\section*{Գործնական 2–3 — Հարաբերություններ}

%% ============================================================
%%  1. Հարաբերություններ
%% ============================================================

\subsection*{Խնդիր 1.1}

\textbf{(a)} 12-ի բաժանարարները $\mathbb{N}$-ում՝
\[
  \boxed{\{(1,12),(2,12),(3,12),(4,12),(6,12),(12,12)\}}
\]

\textbf{(b)} Զույգեր $(d,n)$, որտեղ $d,n\in\{2,3,4,5,6\}$ և $d\mid n$՝
\[
  \boxed{\{(2,2),(2,4),(2,6),(3,3),(3,6),(4,4),(5,5),(6,6)\}}
\]

\textbf{(c)} Եռյակներ $(x,y,z)$, որտեղ $x=y+z$, և բոլորը պատկանում են $\{1,2,3\}$ բազմությանը՝
\[
  \boxed{\{(2,1,1),(3,1,2),(3,2,1)\}}
\]

%% ────────────────────────────────────────

\subsection*{Խնդիր 1.2}

\textbf{(a)} $\boxed{D=\{a,c,d\},\quad R=\{1,2,4,5\}}$։

\textbf{(b)} Օրինաչափություն $(n,2n)$՝ $\boxed{D=\mathbb{N},\quad R=2\mathbb{N}}$։

\textbf{(c)} $x=y^2$՝ $\boxed{D=[0,+\infty),\quad R=\mathbb{R}}$։

\textbf{(d)} $x^2+y^2\le 16$՝ $\boxed{D=[-4,4],\quad R=[-4,4]}$։

\textbf{(e)} $(x,y)\in\mathbb{N}^2,\; x\le y$՝ $\boxed{D=\mathbb{N},\quad R=\mathbb{N}}$։

\textbf{(f)} $(x,y)\in\mathbb{N}^2,\; y\le x$՝ $\boxed{D=\mathbb{N},\quad R=\mathbb{N}}$։

\textbf{(g)} $x+y\ge 0$ $\mathbb{R}^2$-ում՝ $\boxed{D=\mathbb{R},\quad R=\mathbb{R}}$։

\textbf{(h)} $2x\ge 3y$ $\mathbb{R}^2$-ում՝ $\boxed{D=\mathbb{R},\quad R=\mathbb{R}}$։

%% ============================================================
%%  2. Հարաբերությունների ներկայացում
%% ============================================================

\subsection*{Խնդիր 2.1}

$A=\{a,b,c,d,e\}$ բազմության վրա։

\textbf{(a)} Սիմետրիկ է; չուղղորդված կողեր՝ $a$--$b$, $a$--$c$, $b$--$c$, $d$--$e$։
\[
  M = \begin{pmatrix}
  0&1&1&0&0\\1&0&1&0&0\\1&1&0&0&0\\0&0&0&0&1\\0&0&0&1&0
  \end{pmatrix}
\]

\textbf{(b)} Չուղղորդված կողեր՝ $a$--$b$, $a$--$c$, $b$--$c$, $c$--$d$. $e$ գագաթը մեկուսացված է։
\[
  M = \begin{pmatrix}
  0&1&1&0&0\\1&0&1&0&0\\1&1&0&1&0\\0&0&1&0&0\\0&0&0&0&0
  \end{pmatrix}
\]

\textbf{(c)} Նույնական է (b) մասին; նույն գրաֆը և մատրիցը։

%% ────────────────────────────────────────

\subsection*{Խնդիր 2.2}

$A=\{a,b,c,d,e,f\}$ բազմության վրա։ Տողեր/սյունակներ $a,b,c,d,e,f$ հերթականությամբ։

\textbf{(a)} Ուղղորդված կողեր՝ $a\!\to\!b,\; b\!\to\!c,\; a\!\to\!c,\; b\!\to\!e,\; c\!\to\!f,\; c\!\to\!d,\; d\!\to\!f,\; f\!\to\!e$։
\[
  M = \begin{pmatrix}
  0&1&1&0&0&0\\0&0&1&0&1&0\\0&0&0&1&0&1\\0&0&0&0&0&1\\0&0&0&0&0&0\\0&0&0&0&1&0
  \end{pmatrix}
\]

\textbf{(b)} Ուղղորդված կողեր՝ $a\!\to\!b,\; b\!\to\!c,\; d\!\to\!c,\; d\!\to\!e,\; f\!\to\!e,\; f\!\to\!a,\; b\!\to\!e$։
\[
  M = \begin{pmatrix}
  0&1&0&0&0&0\\0&0&1&0&1&0\\0&0&0&0&0&0\\0&0&1&0&1&0\\0&0&0&0&0&0\\1&0&0&0&1&0
  \end{pmatrix}
\]

\textbf{(c)} Ուղղորդված կողեր՝ $a\!\to\!c,\; c\!\to\!c,\; c\!\to\!d,\; c\!\to\!e,\; e\!\to\!c,\; e\!\to\!e,\; e\!\to\!f,\; a\!\to\!e$. $b$ գագաթը մեկուսացված է։
\[
  M = \begin{pmatrix}
  0&0&1&0&1&0\\0&0&0&0&0&0\\0&0&1&1&1&0\\0&0&0&0&0&0\\0&0&1&0&1&1\\0&0&0&0&0&0
  \end{pmatrix}
\]

%% ============================================================
%%  3. Գործողություններ հարաբերությունների հետ
%% ============================================================

\subsection*{Խնդիր 3.1}

\emph{Ապացույցի ուրվագիծ} (${}^{-1}$ և $\circ$ ստանդարտ հատկություններ)։

(i) $(\rho^{-1})^{-1}=\rho$: $(x,y)\in(\rho^{-1})^{-1}\Leftrightarrow (y,x)\in\rho^{-1}\Leftrightarrow (x,y)\in\rho$։

(ii) $(\rho\circ\sigma)^{-1}=\sigma^{-1}\circ\rho^{-1}$: $(z,x)\in(\rho\circ\sigma)^{-1} \Leftrightarrow \exists y\!:\,(x,y)\in\rho\land(y,z)\in\sigma \Leftrightarrow \exists y\!:\,(z,y)\in\sigma^{-1}\land(y,x)\in\rho^{-1} \Leftrightarrow (z,x)\in\sigma^{-1}\circ\rho^{-1}$։ \qed

%% ────────────────────────────────────────

\subsection*{Խնդիր 3.2}

$\rho=\{(x,y)\mid y=x^2+5\}$, $\sigma=\{(x,y)\mid y=3x\}$ $\mathbb{R}$-ի վրա։
Կոմպոզիցիան կարդացվում է ձախից աջ ($\rho\cdot\sigma$-ում սկզբում կիրառվում է $\rho$-ն)։

$\rho\cdot\sigma$: $x\xrightarrow{\rho}x^2+5\xrightarrow{\sigma}3(x^2+5)$։ \;$\boxed{\rho\cdot\sigma=\{(x,\,3x^2+15)\}}$։

$\sigma\cdot\rho$: $x\xrightarrow{\sigma}3x\xrightarrow{\rho}(3x)^2+5$։ \;$\boxed{\sigma\cdot\rho=\{(x,\,9x^2+5)\}}$։

$\boxed{\rho^{-1}=\{(x,y)\mid x=y^2+5\}}$ (այսինքն՝ $y=\pm\sqrt{x-5}$, որոշված է $x\ge 5$-ի համար)։

$\boxed{\sigma^{-1}=\{(x,y)\mid y=x/3\}}$։

$\rho\cdot\sigma^{-1}$: $x\xrightarrow{\rho}x^2+5\xrightarrow{\sigma^{-1}}\tfrac{x^2+5}{3}$։ \;$\boxed{\rho\cdot\sigma^{-1}=\left\{\left(x,\,\tfrac{x^2+5}{3}\right)\right\}}$։

$\rho^{-1}\cdot\sigma$: $x\xrightarrow{\rho^{-1}}\pm\sqrt{x-5}\xrightarrow{\sigma}\pm 3\sqrt{x-5}$ \;($x\ge 5$)։ \;$\boxed{\rho^{-1}\cdot\sigma=\{(x,\,\pm 3\sqrt{x-5})\mid x\ge 5\}}$։

%% ────────────────────────────────────────

\subsection*{Խնդիր 3.3}

$\rho=\{(x,x+2)\}$, $\sigma=\{(x,x^2)\}$ $\mathbb{Z}$-ի վրա։

$\rho\cdot\sigma$: $x\xrightarrow{\rho}x+2\xrightarrow{\sigma}(x+2)^2$։ \;$\boxed{\rho\cdot\sigma=\{(x,\,(x+2)^2)\}}$։

$\sigma\cdot\rho$: $x\xrightarrow{\sigma}x^2\xrightarrow{\rho}x^2+2$։ \;$\boxed{\sigma\cdot\rho=\{(x,\,x^2+2)\}}$։

Դրանք \emph{հավասար չեն} ($x=1$-ում՝ $\rho\cdot\sigma$-ն տալիս է $9$, $\sigma\cdot\rho$-ն՝ $3$)։

%% ────────────────────────────────────────

\subsection*{Խնդիր 3.4}

$\rho=\{(x,y)\in\mathbb{Z}^2\mid x+y \text{ կենտ է}\}$։

$\rho^2$-ու համար՝ պետք է գտնել $z$, որ $x+z$ կենտ է և $z+y$ կենտ է, ուստի $x\equiv y\pmod{2}$։
$\boxed{\rho^2=\{(x,y)\in\mathbb{Z}^2\mid x+y \text{ զույգ է}\}}$։

$\boxed{\rho^3=\rho}$, և ընդհանուր առմամբ $\rho^{2k}=\rho^2$, $\rho^{2k+1}=\rho$։

%% ────────────────────────────────────────

\subsection*{Խնդիր 3.5}

$\rho=\{(x,x+1)\mid x\in\mathbb{N}\}$։

\[\boxed{\rho^k=\{(x,\,x+k)\mid x\in\mathbb{N}\},\quad k\ge 1.}\]

%% ────────────────────────────────────────

\subsection*{Խնդիր 3.6}

$\rho=\{(1,2),(2,3),(3,4),(4,5)\}$։

$\boxed{\rho^2=\{(1,3),(2,4),(3,5)\}}$։

$\boxed{\rho^3=\{(1,4),(2,5)\}}$։

$\boxed{\rho^4=\{(1,5)\}}$։

$\boxed{\rho^5=\varnothing}$։

%% ============================================================
%%  4. Հարաբերությունների հատկությունները
%% ============================================================

\subsection*{Խնդիր 4.1}

\begin{tabular}{l|cccc}
& Ռեֆլ. & Սիմ. & Հակասիմ. & Տրանզ. \\\hline
(a) $x^2+y^2=1$, $\mathbb{R}$ & Ոչ & Այո & Ոչ & Ոչ \\
(b) $x^2=y^2$, $\mathbb{R}$   & Այո & Այո & Ոչ & Այո \\
(c) $x/y$, $\mathbb{Z}$       & Այո & Ոչ  & Ոչ  & Այո \\
(d) $x/y$, $\mathbb{N}$       & Այո & Ոչ  & Այո & Այո \\
(e) $\gcd=1$, $\mathbb{N}$    & Ոչ  & Այո & Ոչ  & Ոչ \\
(f) $\gcd=1$, $\mathbb{Z}$    & Ոչ  & Այո & Ոչ  & Ոչ
\end{tabular}

\smallskip
Այստեղ $x/y$-ն նշանակում է բաժանելիություն ($y\mid x$)։ (b)-ն համարժեքության հարաբերություն է։ (d)-ն (բաժանելիությունը $\mathbb{N}$-ի վրա) մասնակի կարգ է, բայց \emph{ոչ} գծային — օրինակ՝ $2$-ը և $3$-ը անհամեմատելի են։ (c)-ն (բաժանելիությունը $\mathbb{Z}$-ի վրա) միայն նախակարգ է (ռեֆլեքսիվ ու տրանզիտիվ, բայց հակասիմետրիկ չէ), քանի որ $2\mid -2$ և $-2\mid 2$, բայց $2\ne -2$։

%% ────────────────────────────────────────

\subsection*{Խնդիր 4.2}

$\rho=\{(x,y)\in\mathbb{Z}^2\mid |x-y|\text{ զույգ է}\}$։

Ռեֆլեքսիվ — Այո; Սիմետրիկ — Այո; Հակասիմետրիկ — Ոչ; Տրանզիտիվ — Այո։

$\boxed{\text{Համարժեքության հարաբերություն։ Դասեր՝ զույգ ամբողջ թվեր և կենտ ամբողջ թվեր։}}$

%% ────────────────────────────────────────

\subsection*{Խնդիր 4.3}

\textbf{Ռեֆլեքսիվության} պահպանումը ($R,S$ ռեֆլեքսիվ են $A$-ի վրա)՝

(a) $R\cap S$: \textbf{Ճիշտ է։} \quad
(b) $R\cup S$: \textbf{Ճիշտ է։} \quad
(c) $R\circ S$: \textbf{Ճիշտ է։}

(d) $R\setminus S$: \textbf{Սխալ է} ($(x,x)\in R\cap S$, ուստի $(x,x)\notin R\setminus S$)։ \quad
(e) $R\oplus S$: \textbf{Սխալ է} (նույն պատճառով)։

%% ────────────────────────────────────────

\subsection*{Խնդիր 4.4}

\textbf{Սիմետրիկության} պահպանումը ($R,S$ սիմետրիկ են)՝

(a) $R\cap S$: \textbf{Ճիշտ է։} \quad
(b) $R\cup S$: \textbf{Ճիշտ է։}

(c) $R\circ S$: \textbf{Սխալ է։} ($R=\{(1,2),(2,1)\}$, $S=\{(2,3),(3,2)\}$: $(1,3)\in R\circ S$, բայց $(3,1)\notin R\circ S$։)

(d) $R\setminus S$: \textbf{Ճիշտ է։} \quad
(e) $R\oplus S$: \textbf{Ճիշտ է։}

%% ────────────────────────────────────────

\subsection*{Խնդիր 4.5}

\textbf{Հակասիմետրիկության} պահպանումը ($R,S$ հակասիմետրիկ են)՝

(a) $R\cap S$: \textbf{Ճիշտ է։} \quad
(b) $R\cup S$: \textbf{Սխալ է} ($R=\{(1,2)\}$, $S=\{(2,1)\}$):

(c) $R\circ S$: \textbf{Սխալ է։} \quad
(d) $R\setminus S$: \textbf{Ճիշտ է} ($R\setminus S\subseteq R$)։

(e) $R\oplus S$: \textbf{Սխալ է} (նույն հակաօրինակը, ինչ (b)-ում)։

%% ────────────────────────────────────────

\subsection*{Խնդիր 4.6}

\textbf{Տրանզիտիվության} պահպանումը ($R,S$ տրանզիտիվ են)՝

(a) $R\cap S$: \textbf{Ճիշտ է։}

(b) $R\cup S$: \textbf{Սխալ է} ($R=\{(1,2)\}$, $S=\{(2,3)\}$; $(1,3)\notin R\cup S$)։

(c) $R\circ S$: \textbf{Սխալ է։}

(d) $R\setminus S$: \textbf{Սխալ է} ($R=\{(1,2),(2,3),(1,3)\}$, $S=\{(1,3)\}$; $R\setminus S=\{(1,2),(2,3)\}$, տրանզիտիվ չէ)։

(e) $R\oplus S$: \textbf{Սխալ է։}

%% ============================================================
%%  5. Կարգի հարաբերություններ
%% ============================================================

\subsection*{Խնդիր 5.1}

Կառուցել մասնակի կարգավորված բազմություն, որն ունի միակ մինիմալ տարր, բայց չունի ամենափոքր տարր։

\emph{Նշում.} Ցանկացած \emph{վերջավոր} մասնակի կարգում միակ մինիմալ տարրը ավտոմատ կերպով ամենափոքրն է, ուստի անհրաժեշտ է անվերջ մասնակի կարգ (անվերջ նվազող շղթայով)։

\[\boxed{A=\mathbb{Z}^-\cup\{m\}},\]
որտեղ $\mathbb{Z}^-=\{\ldots,-3,-2,-1\}$ ունի սովորական կարգը, իսկ $m$-ն անհամեմատելի է $\mathbb{Z}^-$-ի յուրաքանչյուր տարրի հետ։

Այդ դեպքում $m$-ը \emph{միակ մինիմալ} տարրն է՝ ոչ մի $n\in\mathbb{Z}^-$ մինիմալ չէ (շղթան $\cdots<-2<-1$ անվերջ նվազում է, ուստի յուրաքանչյուր $n$-ի համար $n-1<n$ գտնվում է նրանից ներքև), իսկ $m$-ից ներքև ոչինչ չկա։ Բայց $m$-ը \emph{ամենափոքր} տարր չէ, քանի որ $m\not\le n$ որևէ $n\in\mathbb{Z}^-$-ի համար ($m$-ն անհամեմատելի է նրանց հետ)։ \qed

%% ────────────────────────────────────────

\subsection*{Խնդիր 5.2}

\emph{Ապացույցի ուրվագիծ.} Դիցուք $\rho$-ն մասնակի կարգ է $A$-ի վրա։ Ցույց տանք, որ $\rho^{-1}$-ը ռեֆլեքսիվ, հակասիմետրիկ և տրանզիտիվ է։
Ռեֆլեքսիվ՝ $(x,x)\in\rho\Rightarrow(x,x)\in\rho^{-1}$։
Հակասիմետրիկ՝ $(x,y),(y,x)\in\rho^{-1}\Rightarrow(y,x),(x,y)\in\rho\Rightarrow x=y$։
Տրանզիտիվ՝ $(x,y),(y,z)\in\rho^{-1}\Rightarrow(z,y),(y,x)\in\rho\Rightarrow(z,x)\in\rho\Rightarrow(x,z)\in\rho^{-1}$։ \qed

%% ────────────────────────────────────────

\subsection*{Խնդիր 5.3}

\emph{Ապացույցի ուրվագիծ.} Դիցուք $\{\rho_i\}$-ն մասնակի կարգեր են $A$-ի վրա, $\rho=\bigcap_i\rho_i$։
Ռեֆլեքսիվ՝ $(x,x)\in\rho_i$ բոլոր $i$-երի համար։
Հակասիմետրիկ՝ $(x,y),(y,x)\in\rho\subseteq\rho_i\Rightarrow x=y$։
Տրանզիտիվ՝ $(x,y),(y,z)\in\rho\subseteq\rho_i\Rightarrow(x,z)\in\rho_i$ բոլոր $i$-երի համար $\Rightarrow(x,z)\in\rho$։ \qed

%% ────────────────────────────────────────

\subsection*{Խնդիր 5.4}

\emph{Ապացույցի ուրվագիծ.} $\rho_A=\rho\cap(A\times A)$-ն ժառանգում է ռեֆլեքսիվությունը ($x\in A$-ի համար), հակասիմետրիկությունը և տրանզիտիվությունը $\rho$-ից, քանի որ $\rho_A\subseteq\rho$ և դիտարկվող բոլոր տարրերը պատկանում են $A$-ին։ \qed

%% ────────────────────────────────────────

\subsection*{Խնդիր 5.5}

$\boxed{\rho_1\cup\rho_2\text{-ը գծային կարգ է այն և միայն այն դեպքում, եթե }\rho_1\subseteq\rho_2\text{ կամ }\rho_2\subseteq\rho_1.}$

Եթե մեկը պարունակում է մյուսին, միավորումը հավասար է մեծագույնին։ Հակառակ դեպքում, եթե ոչ մեկը մյուսին չի պարունակում, միավորման հակասիմետրիկությունը կարող է խախտվել։

%% ────────────────────────────────────────

\subsection*{Խնդիր 5.6}

\emph{Ապացույցի ուրվագիծ.} Ինդուկցիա ըստ $|A|=n$-ի։ Հենք՝ $n=1$: ակնհայտ է։ Քայլ՝ վերցնենք $a\in A$; ըստ վարկածի $A\setminus\{a\}$-ն ունի գծային կարգ $a_1<\cdots<a_{n-1}$։ Ավելացնենք $a$-ն վերջում՝ $a_1<\cdots<a_{n-1}<a$։ Սա գծային կարգ է $A$-ի վրա։ \qed

%% ============================================================
%%  6. Համարժեքության հարաբերություններ
%% ============================================================

\subsection*{Խնդիր 6.1}

\textbf{(a)} $xy>0$ $\mathbb{Q}$-ի վրա՝ \textbf{Ոչ։} Ռեֆլեքսիվ չէ ($0\cdot 0=0\not>0$)։

\textbf{(b)} $x-y\in\mathbb{Z}$ $\mathbb{Q}$-ի վրա՝ \textbf{Այո։} Ռեֆլեքսիվ ($x-x=0\in\mathbb{Z}$), սիմետրիկ և տրանզիտիվ ($x-y,\,y-z\in\mathbb{Z}\Rightarrow x-z\in\mathbb{Z}$)։ Դասերը $x+\mathbb{Z}$ հարակից դասերն են՝ միևնույն կոտորակային մասն ունեցող ռացիոնալ թվերը։

\textbf{(c)} $(a\le b)\lor(a>b)$ $\mathbb{N}$-ի վրա՝ \textbf{Այո։} Քանի որ ցանկացած $a,b$-ի համար $a\le b$ կամ $a>b$՝ պայմանը միշտ ճիշտ է, ուստի $\rho=\mathbb{N}\times\mathbb{N}$-ը համընդհանուր հարաբերությունն է — տրիվիալ կերպով ռեֆլեքսիվ, սիմետրիկ, տրանզիտիվ (մեկ դաս՝ ամբողջ $\mathbb{N}$-ը)։

\textbf{(d)} $|a-b|\le 2$ $\mathbb{N}$-ի վրա՝ \textbf{Ոչ։} Տրանզիտիվ չէ ($|1-3|=2\le 2$, $|3-5|=2\le 2$, բայց $|1-5|=4>2$)։

%% ────────────────────────────────────────

\subsection*{Խնդիր 6.2}

(Համարժեքության հարաբերություն $\Leftrightarrow$ տրոհում։)

\emph{Ապացույցի ուրվագիծ.}
$(\Rightarrow)$ Դասերը ծածկում են $A$-ն (ռեֆլեքսիվություն՝ $x\in[x]$): Եթե $[a]\cap[b]\ne\varnothing$, վերցրեք $c$ երկուսից էլ; ապա $a\rho c$ և $c\rho b$, ուստի $a\rho b$ և $[a]=[b]$։
$(\Leftarrow)$ Սահմանեք $x\rho y$ այն և միայն այն դեպքում, եթե $x,y$-ը նույն բլոկում են; ռեֆլեքսիվությունը, սիմետրիկությունը, տրանզիտիվությունը անմիջական են։ \qed

%% ────────────────────────────────────────

\subsection*{Խնդիր 6.3}

\textbf{$\mathbb{Z}_2$:}
\[
\begin{array}{c|cc}
+ & [0] & [1] \\\hline
{[0]} & [0] & [1] \\
{[1]} & [1] & [0]
\end{array}
\qquad
\begin{array}{c|cc}
\cdot & [0] & [1] \\\hline
{[0]} & [0] & [0] \\
{[1]} & [0] & [1]
\end{array}
\]

\textbf{$\mathbb{Z}_3$:}
\[
\begin{array}{c|ccc}
+ & [0] & [1] & [2] \\\hline
{[0]} & [0] & [1] & [2] \\
{[1]} & [1] & [2] & [0] \\
{[2]} & [2] & [0] & [1]
\end{array}
\qquad
\begin{array}{c|ccc}
\cdot & [0] & [1] & [2] \\\hline
{[0]} & [0] & [0] & [0] \\
{[1]} & [0] & [1] & [2] \\
{[2]} & [0] & [2] & [1]
\end{array}
\]

\textbf{$\mathbb{Z}_4$:}
\[
\begin{array}{c|cccc}
+ & [0] & [1] & [2] & [3] \\\hline
{[0]} & [0] & [1] & [2] & [3] \\
{[1]} & [1] & [2] & [3] & [0] \\
{[2]} & [2] & [3] & [0] & [1] \\
{[3]} & [3] & [0] & [1] & [2]
\end{array}
\qquad
\begin{array}{c|cccc}
\cdot & [0] & [1] & [2] & [3] \\\hline
{[0]} & [0] & [0] & [0] & [0] \\
{[1]} & [0] & [1] & [2] & [3] \\
{[2]} & [0] & [2] & [0] & [2] \\
{[3]} & [0] & [3] & [2] & [1]
\end{array}
\]

%% ────────────────────────────────────────

\subsection*{Խնդիր 6.4}

$[a]\in\mathbb{Z}_n$ հակադարձելի է այն և միայն այն դեպքում, եթե $\gcd(a,n)=1$։

\emph{Ապացույցի ուրվագիծ.}
$(\Rightarrow)$ $ab\equiv 1\pmod{n}$ նշանակում է $ab - kn = 1$, ուստի $\gcd(a,n)=1$ ըստ Բեզուի։
$(\Leftarrow)$ $\gcd(a,n)=1$ հետևում է $\exists\,b,k$: $ab+kn=1$, ուստի $ab\equiv 1\pmod{n}$։ \qed

%% ============================================================
%%  7. Գործողություններ համարժեքության հարաբերությունների հետ (Փակույթներ)
%% ============================================================

\subsection*{Խնդիր 7.1}

$A=\{a,b,c\}$-ի վրա, $\Delta_A=\{(a,a),(b,b),(c,c)\}$։

\textbf{(a)} $\rho=\{(a,b),(b,a)\}$։
$r(\rho)=\rho\cup\Delta_A=\{(a,a),(b,b),(c,c),(a,b),(b,a)\}$։
$s(\rho)=\rho$ (արդեն սիմետրիկ է)։
$t(\rho)=\{(a,b),(b,a),(a,a),(b,b)\}$։

\textbf{(b)} $\rho=\{(a,b),(b,c)\}$։
$r(\rho)=\{(a,a),(b,b),(c,c),(a,b),(b,c)\}$։
$s(\rho)=\{(a,b),(b,a),(b,c),(c,b)\}$։
$t(\rho)=\{(a,b),(b,c),(a,c)\}$։

\textbf{(c)} $\rho=\{(a,a),(a,b),(c,b),(c,a)\}$։
$r(\rho)=\rho\cup\{(b,b),(c,c)\}$։
$s(\rho)=\rho\cup\{(b,a),(b,c),(a,c)\}$։
$t(\rho)=\rho$ (արդեն տրանզիտիվ է՝ $(c,a)\circ(a,b)=(c,b)\in\rho$)։

\textbf{(d)} $\rho=\{(a,a),(a,c),(b,c)\}$։
$r(\rho)=\rho\cup\{(b,b),(c,c)\}$։
$s(\rho)=\rho\cup\{(c,a),(c,b)\}$։
$t(\rho)=\rho$ (արդեն տրանզիտիվ է)։

\textbf{(e)} $\rho=\{(a,b),(c,a),(b,c)\}$։
$r(\rho)=\rho\cup\Delta_A$։
$s(\rho)=\{(a,b),(b,a),(c,a),(a,c),(b,c),(c,b)\}$։
$t(\rho)=A\times A$ (ամբողջական ցիկլը $a\to b\to c\to a$ գեներացնում է բոլոր 9 զույգերը)։

\textbf{(f)} $\rho=\{(a,b),(c,b),(b,c)\}$։
$r(\rho)=\rho\cup\Delta_A$։
$s(\rho)=\{(a,b),(b,a),(c,b),(b,c)\}$։
$t(\rho)=\{(a,b),(a,c),(b,c),(b,b),(c,b),(c,c)\}$։

%% ────────────────────────────────────────

\subsection*{Խնդիր 7.2}

$A=\{1,2,3,4\}$, $\rho=\{(1,2),(3,1),(3,2),(2,4)\}$։

$\boxed{r(\rho)=\{(1,1),(2,2),(3,3),(4,4),(1,2),(2,4),(3,1),(3,2)\}}$։

$\boxed{s(\rho)=\{(1,2),(2,1),(1,3),(3,1),(2,3),(3,2),(2,4),(4,2)\}}$։

$\boxed{t(\rho)=\{(1,2),(1,4),(2,4),(3,1),(3,2),(3,4)\}}$։

(Նոր զույգեր $\rho^2$-ից՝ $(1,4)$ ճանապարհով $1\!\to\!2\!\to\!4$; $(3,4)$ ճանապարհով $3\!\to\!2\!\to\!4$; $(3,2)$ ճանապարհով $3\!\to\!1\!\to\!2$ արդեն կա: 4-ը մտնող կող չունի, ուստի $\rho^3$-ը ոչինչ չի ավելացնում։)

%% ────────────────────────────────────────

\subsection*{Խնդիր 7.3}

\emph{Ապացույցի ուրվագիծ.} $\tau=\rho\cap\sigma$, որտեղ $\rho,\sigma$ համարժեքության հարաբերություններ են։
Ռեֆլեքսիվ՝ $(x,x)\in\rho\cap\sigma$։
Սիմետրիկ՝ $(x,y)\in\tau\Rightarrow(y,x)\in\rho$ և $(y,x)\in\sigma\Rightarrow(y,x)\in\tau$։
Տրանզիտիվ՝ $(x,y),(y,z)\in\tau\subseteq\rho\Rightarrow(x,z)\in\rho$; նմանապես $\in\sigma$; ուստի $(x,z)\in\tau$։ \qed

%% ────────────────────────────────────────

\subsection*{Խնդիր 7.4}

$\rho\circ\theta$-ն համարժեքության հարաբերություն է այն և միայն այն դեպքում, եթե $\rho\circ\theta=\theta\circ\rho$։

\emph{Ապացույցի ուրվագիծ.}
$(\Rightarrow)$ $(\rho\circ\theta)^{-1}=\theta^{-1}\circ\rho^{-1}=\theta\circ\rho$: Քանի որ $\rho\circ\theta$-ն սիմետրիկ է, $\rho\circ\theta=\theta\circ\rho$։
$(\Leftarrow)$ Ռեֆլեքսիվություն՝ $\Delta\subseteq\rho,\theta\Rightarrow\Delta\subseteq\rho\circ\theta$: Սիմետրիկություն՝ $(\rho\circ\theta)^{-1}=\theta\circ\rho=\rho\circ\theta$: Տրանզիտիվություն՝ $(\rho\circ\theta)^2=\rho\circ(\theta\circ\rho)\circ\theta =\rho\circ(\rho\circ\theta)\circ\theta\subseteq\rho^2\circ\theta^2\subseteq\rho\circ\theta$։ \qed

%% ────────────────────────────────────────

\subsection*{Խնդիր 7.5}

$\rho\cup\theta$-ն համարժեքության հարաբերություն է այն և միայն այն դեպքում, եթե $\rho\circ\theta=\theta\circ\rho$։

\emph{Ապացույցի ուրվագիծ.}
$(\Rightarrow)$ Եթե $\rho\cup\theta$-ն համարժեքություն է (հետևաբար տրանզիտիվ), ապա $\rho\circ\theta\subseteq(\rho\cup\theta)\circ(\rho\cup\theta)\subseteq\rho\cup\theta$: Նմանապես $\theta\circ\rho\subseteq\rho\cup\theta$: Քանի որ $\rho\cup\theta\subseteq\rho\circ\theta$ (քանի որ $\rho=\rho\circ\Delta\subseteq\rho\circ\theta$), ստանում ենք $\rho\cup\theta=\rho\circ\theta$, և սիմետրիայից նաև $=\theta\circ\rho$։

$(\Leftarrow)$ Եթե $\rho\circ\theta=\theta\circ\rho$, ապա ըստ Խնդիր 7.4-ի, $\rho\circ\theta$-ն համարժեքության հարաբերություն է։ Ցույց է տրվում, որ $\rho\cup\theta=\rho\circ\theta$, որը հետևաբար համարժեքության հարաբերություն է։ \qed